%==============================================================================
% Appendix A: Topology conventions, local reductions, and caveats
%==============================================================================
\section{Topology conventions, local reductions, and caveats}
\label{app:A}

\subsection{Structure constants and $\chi$ values}
\label{app:A:structure}

The spatial structure constants for the four topologies are given as
specialisations of the five-parameter LZ-native isotropic family.  The
family parameters are $(a,b,c,u,v)$, and the antisymmetrised structure
constants are defined by
\begin{equation}
  C^{1}{}_{23}=\frac{a}{q},\quad
  C^{2}{}_{31}=\frac{b}{q},\quad
  C^{3}{}_{12}=\frac{c}{q},\quad
  C^{2}{}_{12}=\frac{u}{q},\quad
  C^{3}{}_{13}=\frac{v}{q}.
\end{equation}
The three-dimensional Ricci scalar then reads
\begin{equation}
  {}^{(3)}R\,q^{2} = -\tfrac{1}{2}\,\Disc,
  \qquad
  \Disc = a^{2}-2ab-2ac+b^{2}-2bc+c^{2}+4u^{2}+4uv+4v^{2},
\end{equation}
and the $\chi$ of this paper is defined by
\begin{equation}
  \chi = -\frac{\Disc}{12}.
\end{equation}
The parameter values and $\chi$ values for each topology are listed in
Table~\ref{tab:appA_topology}.

\begin{table}[H]
\centering
\begingroup
\renewcommand{\arraystretch}{2.0}
\begin{tabularx}{\linewidth}{lcrrYY}
\toprule
Topology & $(a,b,c,u,v)$ & $\Disc$ & $\chi$ & Reduction status & Caveat \\
\midrule
$\Sthree$   & $(4,4,4,0,0)$  & $-48$ & $4$     & closed isotropic reference   & --- \\
$\Tthree$   & $(0,0,0,0,0)$  & $0$   & $0$     & flat reference               & --- \\
$\Nilthree$ & $(0,0,-1,0,0)$ & $1$   & $-1/12$ & isotropic-scale local branch & anisotropic EOM left to future work \\
$\Solthree$ & $(0,0,0,1,-1)$ & $4$   & $-1/3$  & isotropic-scale local branch & compact quotient outside the scope of this paper \\
\bottomrule
\end{tabularx}
\endgroup
\caption{Topology conventions, $\chi$ values, and reduction status.}
\label{tab:appA_topology}
\end{table}

\paragraph{Verification.}
$\Sthree$: $\Disc=(16-32-32+16-32+16)=-48$, $\chi=4$.
$\Nilthree$: $\Disc=1$, $\chi=-1/12$.
$\Solthree$: $\Disc=(4-4+4)=4$, $\chi=-1/3$.

\subsection{Distinguished loci in the five-parameter family}
\label{app:A:loci}

Table~\ref{tab:appA_loci} records, on top of the topology ordering and the
$(\Disc,\chi)$ conventions fixed in Table~\ref{tab:appA_topology}, only the
additional bridge-facing information.  Here $\Ctop$ is computed from
$\Ctop = 3\Disc/4$ of Appendix~\ref{app:B:bridge}, and each row satisfies
$\Ctop+9\chi=0$.

\begin{table}[H]
\centering
\begin{tabular}{lrll}
\toprule
Topology & $\Ctop=3\Disc/4$ & Structural locus & paper~IV echo \\
\midrule
$\Sthree$   & $-36$ & isotropic class-A diagonal      & triaxial \\
$\Tthree$   & $0$   & origin                          & none \\
$\Nilthree$ & $3/4$ & single class-A axis             & uniaxial \\
$\Solthree$ & $3$   & balanced class-B diagonal pair  & biaxial \\
\bottomrule
\end{tabular}
\caption{Bridge-facing structural loci.}
\label{tab:appA_loci}
\end{table}

The paper~IV echo column is interpretive only.  It records a structural
resonance with the response dictionary of paper~IV, not an equivalence of
mechanisms or a shared theorem.

\subsection{Reduction scope for $\Nilthree$ and $\Solthree$}
\label{app:A:reduction}

$\Nilthree$ is the non-Abelian homogeneous geometry determined by the
left-invariant metric on the Heisenberg
group~\cite{Thurston1997,Scott1983}; the sign convention adopted here is
$C^{3}{}_{12}=-1/q$ (that is, $c=-1$).  Under the isotropic-scale reduction
a single scale parameter $q$ is shared across all directions.  The
anisotropic EOM obtained by introducing independent scale factors is not
derived in this paper.

$\Solthree$ is the geometry based on a solvable Lie
group~\cite{Thurston1997,Scott1983} and enters the five-parameter family
as a balanced class-B diagonal pair $u=1$, $v=-1$.  This is the
isotropic-scale reduced entry of the present paper.  A compact realisation
of $\Solthree$ requires a mapping-torus construction; the choice of lattice
normalisation, the classification of spin structures, and the determination
of spectral data are outside the scope of the present paper.  The
$\Solthree$ entries of Table~\ref{tab:four_topology_admissibility} and
Table~\ref{tab:chi_sign_atlas} are restricted to this local / isotropic
reduced branch.

\subsection{Lorentzian EC+NY action conventions}
\label{app:A:action}

We adopt the orthonormal frame in Lorentzian signature $(-,+,+,+)$, and the
Nieh--Yan coupling is parametrised by real $\thetaNY$.  The frame
block-diagonal structure of $\Pform$ is common to the Lorentzian and
Euclidean settings, but the sign convention of the Hamiltonian constraint
follows the Lorentzian convention
(Section~\ref{sec:hamiltonian:constraint}).  The definition of $\Pint$ and
the orientation of $\QNY$ inherit the conventions of
paper~V~\cite{paper5}.

The constant term $\Ctop=-36$ of the $\Sthree$ specialisation agrees with
the constant term of the $\MX$-$P$ expression of paper~III (``Unified
Geometric Landau EFT'')~\cite{paper3}.  On the other hand, the sign of the
$V^{2}q^{2}$ term originates from the difference between the Euclidean
convention of paper~III~\cite{paper3} and the Lorentzian metric-aware
epsilon / legacy temporal-index map of the present paper.
